\honest{Cultish Countercultishness}{Eliezer Yudkowsky}{2007}

Joining a cult is ``probably one of the worse things that can happen to you.'' Real cults use love bombing, sleep deprivation, isolation and confession, and take all your money. So I should sympathize with people who are nervous that they might be joining one. Instead, their nervousness grates on my nerves.

Point one: cults are not a natural kind, like cats and dogs. When group failures feed each other and combine, the result ``is not a cult essence; it is a cult attractor.'' So do not settle whether a group is a cult by finding one item on a checklist; look at each characteristic separately. And a cult can hold a true idea: ``Cultishness is a characteristic of groups more than hypotheses.'' \nb{The attractor, the entropy image and the claim that every cause wants to be a cult are from Yudkowsky's ``Every Cause Wants to Be a Cult,'' posted eighteen days earlier.}

The second error is that the nervous asker wants reassurance, not an answer, and once reassured stops pushing against the slide into cultishness that every cause must resist. \nb{This is the lesson of the first of Yudkowsky's ``Two Cult Koans,'' posted nine days earlier.}

The third mistake, ``I strongly suspect,'' is that the nervousness comes from the wrong place. People ask it of cryonics but not of political rallies, where the same group errors occur. ``So why be any more worried about having your head frozen after you stop breathing?'' \nb{A rally costs an evening. Cryonics cost from \$28,000 to \$200,000 in the United States as of 2011, for storage in the hope that revival may one day be possible. The post does not mention the difference.} I suspect the nervousness is the fear of lonely dissent, the feeling of Asch's subjects facing a unanimous room. The next paragraph says ``That's why'' old religions and parties do not make people nervous, and after that the nervousness simply is ``the fear of believing something that will make your friends look at you really oddly.'' I add no evidence between the suspicion and the conclusion.

I grant that ``socially unusual belief puts a group at risk for ingroup-outgroup thinking and evaporative cooling.'' It is ``a risk factor, not a disease.''

Problem four: cults exploit the fear of lonely dissent, and the wish to be sure. Conformity ``glues you to wherever you are.'' Problem five: the right attitude is caution about what a group may become. I share that caution about AI ideas, and try to place my Go stones in advance.

I am asked ``Are you sure this isn't a cult?'' more often about powerful AIs than about probability theory. ``I don't know if one risk factor is higher than the other, but I know which one sounds weirder.'' \nb{By the post's own account, sounding weirder is a risk factor, and the post has just named another that is specific to the AI cause, the ``Super Happy Agent.''}

Problem six: answering the question at all risks giving the reassurance an ``Evil Guru'' would give, and I feel this essay may do so. It is ``a wrong question.'' Look at the group's reasoning yourself. ``When you know what you're trying to do and why, you'll know whether you're getting it done or not, and whether a group is helping you or hindering you.'' \nb{The essay had said that ``having a nice goal always puts you at risk of the happy death spiral.'' Knowing one's goal does not show whether a group serves it.}

A postscript adds a seventh annoyance: all of this takes a long time to explain.
